\begin{helper}
Let \(U^-_x\) and \(U^+_x\) be first-witness pieces formed from residual
witness predicates
\[
R^-_x=C^-_x\setminus D^-,
\qquad
R^+_x=C^+_x\setminus D^+,
\]
where \(D^-\) and \(D^+\) are global boundary events, not events indexed by
the selected witness \(x\).  Then the selected first-witness pieces are
exactly
\[
U^-_x=\{s:s\in C^-_x,\ s\notin D^-,\text{ and no earlier }z<x
\text{ has }s\in C^-_z\setminus D^-\},
\]
and
\[
U^+_x=\{t:t\in C^+_x,\ t\notin D^+,\text{ and no earlier }z<x
\text{ has }t\in C^+_z\setminus D^+\}.
\]
In particular, an outside-\(B_k\) selected witness does not remove the global
boundary complements; those complements remain part of the residual surface.
\end{helper}
\begin{proof}
Unfold the definition of a residual witness.  On the source side,
\(s\in R^-_x\) means exactly \(s\in C^-_x\) and \(s\notin D^-\), because
\(R^-_x=C^-_x\setminus D^-\).  The first-witness piece \(U^-_x\) adds the
condition that for every earlier index \(z<x\), the same state \(s\) is not
in \(R^-_z\).  Replacing \(R^-_z\) by its definition gives precisely
\(\neg(s\in C^-_z\setminus D^-)\), equivalently
\(\neg(s\in C^-_z\ \text{and}\ s\notin D^-)\).  This proves the displayed
source equality.

The target equality is identical, with \(C^+_z,D^+,R^+_z,U^+_z\) in place of
the source objects.  The argument never uses any relation between the
selected index \(x\) and the global boundary events.  Therefore if, in an
application, \(x\notin B_k\), that fact can rule out only \(x\) as a possible
current-boundary witness; it cannot delete the global complements \(D^-\) and
\(D^+\), which may still be witnessed by a different candidate.
\end{proof}
